\begin{abstract}
% High-performance GPU kernels are essential to efficient AI workloads. Agentic optimization has shown promising results in generating and improving such kernels. However, existing agentic workflows primarily target isolated kernels or model inference, leaving end-to-end training performance largely unexplored.
% We present \method, a compiler-guided agentic system for optimizing automatic-differentiation (AD) kernels. 
% Its key insight is to expose forward computation, backward computation, and their saved-state interface jointly to the agent's search space while using training-context verification and evaluation to select deployable kernels. For each kernel target, \method synthesizes editable seeds from compiler-derived semantics, preserves selected library calls, and jointly evolves the paired implementations. To ensure training-specific correctness and performance, \method verifies and evaluates kernels using numerical, latency, and memory signals from training. Across five agentic workflows, only 17 of 28 accepted implementations with deployed-operator speedups retain statistically supported training gains. \method produces accepted combined deployments on five of six evaluated workloads, achieving up to a $1.119\times$ training speedup over eager PyTorch. \haochen{need to update the experiment results}
Recent advances in agentic GPU kernel optimization have shown promising results in automating the development of high-performance kernels. However, existing approaches primarily target forward computation and inference, and develop and evaluate kernels in isolation from their target execution frameworks. This disconnect is particularly problematic for AI training, where forward and backward computations share intermediate state, numerical discrepancies can accumulate across training steps, and local kernel speedups may not translate into training benefits.
We present \system, an end-to-end agentic system for developing and deploying GPU training kernels. \system integrates agentic kernel development with PyTorch's compiler and runtime. It synthesizes editable implementations from compiler-derived AD semantics, jointly optimizes forward computation, backward computation, and their saved-state interface, and verifies and evaluates generated kernels within their target training workloads.
We also introduce \system-Bench, a training-integrated benchmark comprising 149 operator configurations from nine AI models. Our study across five agentic kernel development workflows reveals that isolated kernel assessment overestimates both correctness and performance. Only 17 of 28 accepted implementations with operator-level speedups achieve statistically supported training gains. ADopt [insert final optimization effectiveness, deployment coverage, latency, and memory results].
\end{abstract}


% \keywords{Triton, GPU kernels, automatic differentiation, code generation,
% program optimization}